\section{Coefficient collisions and higher-dimensional anticodes}
\label{sec:lattice}
\newcommand{\Z}{\mathbb Z}\newcommand{\R}{\mathbb R}

In the GRH project's construction of admissible data, the lemma
\texttt{exists\_three\_term\_free} chooses integer exponents $e$ modulo $o$.
For positive $a,b$ with $a+b\le L$, the congruence
$(a+b)e_2=ae_1+be_3$ forces $e_1=e_2=e_3$.
Its proof uses carry-free base expansions and strict convexity on a sphere.
The source identifies this step as a Behrend-type construction of known
provenance [1]. It suggests changing which coefficient relations must be
excluded, and asking how small an encoding group suffices.

Allow every equality between two sums of $h$ labels, with repetitions,
as a potential collision. A label set free of such collisions between distinct
multisets is a classical $B_h$ set. A coefficient difference
$z=(z_0,\ldots,z_n)$ has zero total and
$\tfrac12\sum_i|z_i|\le h$ after common labels cancel. Conversely, each
such difference comes from two $h$-term sums by padding both with the same
label. Thus the coefficient geometry is exactly that of the root lattice
\[
 A_n=\{z\in\Z^{n+1}:\textstyle\sum_i z_i=0\},\qquad
 d(z,w)=\tfrac12\textstyle\sum_i|z_i-w_i|.
\]
This classical correspondence [2] changes the GRH construction's modular
condition in a specific way. The consequences below are independent of the
analytic and automorphic theorems in that project.

\subsection{Exact higher-dimensional cases}
Write $a_n(D)$ for the largest cardinality of a subset of $A_n$ of diameter at
most $D$. Kova\v cevi\'c [3, \S7] conjectures that the balanced difference
of two discrete simplices attains $a_n(D)$ for every $n,D$. His theorem proves
this for $n=3$ and every $D$. The following are computer-assisted exact instances.
\begin{center}
\begin{tabular}{@{}r*{9}{r}@{}}\toprule
$D$ & 2 & 3 & 4 & 5 & 6 & 7 & 8 & 9 & 10\\\midrule
$a_4(D)$ & 21 & 55 & 131 & 255 & 471 & 780 & 1251 & 1875 & 2751\\
$a_5(D)$ & 31 & 96 & 271 & 606 & 1281 & & & &\\\bottomrule
\end{tabular}
\end{center}
For each displayed pair, every maximizing family is a translate of the balanced construction or its negative. The two $D=2$ entries also
follow from classical Singer difference sets for the prime powers $4$ and $5$,
through the code--anticode bound [2,3]. The cited theorem in rank three
does not cover either row for $D\ge3$. Priority for individual table
entries or for the small-group constants below has not been established.

For example, a family of $K$-multisets on five symbols in which every pair
shares at least $K-3$ elements, counting multiplicities, has size at most
$55$. Equality holds exactly when $K\ge6$.


\subsection{Proof of the table and the multiset interpretation}
For $x\in\Z^n$ put $p(x)=\sum_i\max(x_i,0)$ and $q(x)=p(-x)$.
The isometry $x\mapsto(x,-\sum x_i)$ identifies $d$ with the norm
$\max(p,q)$ of the difference. Put $P=\lceil D/2\rceil$,
$Q=\lfloor D/2\rfloor$. The candidate $S_n(P,Q)=\{x:p(x)\le P,\,q(x)\le Q\}$
has diameter at most $P+Q$, since $p(x-y)\le p(x)+q(y)$ and likewise
for $q$. Choosing the $j$ strictly positive coordinates and using
stars and bars on each sign gives its cardinality
\[
 M_n(D)=\sum_{j=0}^{\min(n,P)}
 {n\choose j}{P\choose j}{Q+n-j\choose n-j}.
\]

\textbf{Upper-bound algorithm.}
Let $F$ be an anticode in root coordinates, and put $m=n+1$.
Subtract each coordinate's minimum on $F$. All coordinates now lie
in $[0,D]$, and all sums equal some $k$. Reflection $x_i\mapsto D-x_i$
allows $0\le k\le\lfloor mD/2\rfloor$. Each possible ambient layer
$V_k$ is finite. For $x,y$ in it,
\[
 d(x,y)=\max_{J\subseteq\{1,\ldots,m\}}
 \sum_{i\in J}(x_i-y_i).
\]
Indeed, the sum of all positive coordinate differences is half their total
absolute value. Accordingly $F$ has diameter at most $D$ precisely when
each subset-sum has range at most $D$. Singletons already satisfy this,
and a cut and its complement have equal range.

For each remaining cut let $a,b$ be its minimum and maximum on the whole
layer. If $b-a>D$, try each integer interval $[s,s+D]$ for
$a\le s\le b-D$. Any feasible range is contained in one of these
intervals, by clipping its lower endpoint to $b-D$ when necessary.
If $b-a\le D$, the entire layer satisfies the cut. Intersect with a
chosen interval and recurse. Stop when the intersection has fewer than
$M_n(D)$ points. Every feasible family fits some branch at every depth.
Thus the absence of a larger intersection on all full branches proves the
upper bound. At equality the family must be the entire intersection.
For every leaf of size at least $M_n(D)$ the program checks its size
and compares the full set, after subtracting coordinate minima, with the
two normalized candidate shapes described below. This proves the equality
classification. The supplied \texttt{anticodes.py} performs these
operations with exact integer bit sets. Its outputs give the node counts.

\textbf{Multisets.} For two count vectors of sum $K$, intersection
cardinality is $K-d(x,y)$, so the upper bound holds on every layer.
For the lower bound take $x\in S_n(P,Q)$ and add final coordinate
$P-Q-\sum_i x_i$. All $m$ coordinates are at least $-Q$ and their sum is
$P-Q$. Each coordinate attains $-Q$, so the normalized shape is obtained by
adding $Q$ everywhere. It has coordinate sum $mQ+P-Q$. The negative
shape interchanges $P,Q$. Nonnegative translations give all larger sums.
The equality classification therefore gives the exact threshold
$K\ge mQ+P-Q$, since $P\ge Q$ and $m\ge2$. For $n=4,D=3$ this is $K\ge6$.

\textbf{Approximation in every fixed rank.}
Let $B=\{x\in\R^n:p(x),q(x)\le1\}$. Attaching unit cubes in the
$\Z^n$ coordinates to a family of size $a_n(D)$ gives a union $U$ of
that volume with $U-U\subseteq(D+n)B$. Brunn--Minkowski gives
$a_n(D)\le2^{-n}(D+n)^n\operatorname{vol}(B)$.
The formula for $M_n(D)$ gives leading coefficient
$2^{-n}{2n\choose n}/n!$. This equals $2^{-n}\operatorname{vol}(B)$.
Sum the orthant volumes $\binom nj/(j!(n-j)!)$ over $j$.
Consequently $a_n(D)=M_n(D)+O_n(D^{n-1})$. This elementary
approximation leaves the exact higher-dimensional conjecture open.

\subsection{The least order of a three-sum encoding with five labels}
Let $\phi(3,4)$ be the least order of a finite abelian group containing
a $B_3$ set of cardinality five, in the notation of [2].
The exact minimum is $64$. The minimum among cyclic groups is $65$.
Explicit witnesses in their respective groups are
\[
 \begin{array}{ll}
 (\Z/8\Z)^2 & (0,0),(0,1),(1,0),(4,6),(7,6),\\
 \Z/65\Z & 0,1,4,15,24.
 \end{array}
\]
Each witness gives $\binom{7}{3}=35$ distinct unordered three-label sums.

\textbf{Lower bound.} Translate a candidate so that $b_0=0$. The group
homomorphism $A_4\to G$ defined by $z\mapsto\sum_{i=0}^4z_i b_i$ has
no nonzero kernel element of distance at most $3$, by cancellation and
padding. It is injective on the $55$ points of $S_4(2,1)$ in root
coordinates. Thus $|G|\ge55$. For orders $55$ through $63$, the structure
theorem for finite abelian groups gives the cyclic group at every order
and just four additional groups with invariant factors
$(2,28),(2,2,14),(2,30),(3,21)$.

The check in \texttt{sidon.py} orders each group lexicographically and
enumerates all increasing four-label extensions of $\{0\}$. It retains
a partial set $T$ exactly when its unordered triple sums are distinct.
When adjoining a label $x$, the new sums are the elements of
$(x+2T)\cup(2x+T)\cup\{3x\}$, where $2T$ denotes the pairwise sumset.
There must be precisely $|2T|+|T|+1$ new values, all outside $3T$.
The cardinality test counts all new multisets correctly because the
partial $B_3$ property implies pairwise sums are distinct
by adding one fixed old label. Removing a label preserves the property,
so every valid set is reached. All thirteen groups below order $64$
in the specified range fail. The cyclic group of order $64$ also fails.
Both displayed witnesses pass an additional direct check over the $35$ multisets.

In particular, a full-rank sublattice of $A_4$ of minimum distance at
least four has index at least $64$. The five coordinate labels in its
quotient (translate a common coordinate vector to zero) have the $B_3$
property. A diameter-perfect lattice code at $D=3$ would have index
$a_4(3)=55$. This rules it out.

\textbf{Reproduction.}
From the directory \path{checks/lattice/} run
\begin{verbatim}
python anticodes.py 4 2 3 4 5 6 7 8 9 10
python anticodes.py 5 2 3 4 5 6
python sidon.py
\end{verbatim}

\paragraph{Sources.} [1] L. Sernau et al., \emph{GRH for Dirichlet
L-functions, in Lean 4}, readable edition,
\href{https://github.com/lsernau/grh-lean/blob/72fead102f663d5c7a407964a13dc0cf9d977207/grh/src/New/Reduction/Data/Existence.lean\#L57-L110}{\texttt{Existence.lean}, lines 57--110}.
[2] M. Kova\v cevi\'c and V.~Y.~F. Tan, \emph{Improved Bounds on
Sidon Sets via Lattice Packings of Simplices}, SIAM J. Discrete Math. 31
(2017), 2269--2278, \href{https://doi.org/10.1137/16M1099182}{DOI}.
[3] M. Kova\v cevi\'c, \emph{An Isodiametric Theorem and Lattice
Diameter-Perfect Codes in $A_3$}, 23 July 2026,
\href{https://arxiv.org/abs/2607.21037}{arXiv:2607.21037}, especially
\S\S2, 6--7.
